On prépare une solution aqueuse $\mathrm{(S_1)}$ de sulfate de fer(II)
($\ce{Fe^2+_{(aq)} + SO^2-_{4(aq)}}$) de couleur \textbf{jaune} en dissolvant
une masse $m$ de sulfate de fer(II) $\ce{FeSO4_{(s)}}$ dans un volume
$V=\qty{0.5}{\L}$ de l'eau distillée.
\begin{enumerate}
    \item Donner l'équation de dissolution de $\ce{FeSO4_{(s)}}$ dans l'eau.
    \item Soit $C_1$ la concentration molaire de la solution $\mathrm{(S_1)}$,
          exprimer $C_1$ en fonction de $m$, $M$ et $V$.
\end{enumerate}

On dose un volume $V=\qty{20}{\mL}$ de la solution $\mathrm{(S_1)}$ (avec
quelques goutes d'acide sulfurique) et on lui ajoute progressivement une
solution aqueuse $\mathrm{(S_2)}$ de dichromate de potassium
$(\ce{2K^{+}_{(aq)} + Cr2O^{2-}_{7(aq)}})$ de couleur orange avec une
concentration $C_2=\qty{5e-2}{\mol\per\L}$.

% Les $\ce{Cr2O^{2-}_{7(aq)}}$ ions se réduisent en ions $\ce{Cr^{3+}_{(aq)}}$ et
% les ions $\ce{Fe^{2+}_{aq}}$ s'oxydent en ions $\ce{Fe^{3+}_{(aq)}}$.

L'équivalence est obtenue lorsque le volume ajouté de la solution
$\mathrm{(S_2)}$ est $V_{2,\mathrm{E}}=\qty{10}{\mL}$.
\begin{enumerate}[resume]
    \item Déterminer l'ion responsable de la couleur de chacune de ses deux
          solutions.
    \item Donner le dispositif expérimental permettant de faire le dosage
          direct.
    \item Donner les couples rédox mis en jeux.
    \item Écrire les demi-équations rédox de chaque couple.
    \item Déduire l'équation de la réaction de dosage.
    \item Construire le tableau d'avancement de la réaction du dosage.
    \item Déduire la relation d'équivalence.
    \item Calculer la concentration de la solution $\mathrm{(S_1)}$, et déduire
          la masse $m$ de sulfate de fer(II) $\ce{FeSO4_{(s)}}$.

          On donne~: $M(\ce{Fe})=\qty{56}{\g\per\mol}$~;
          $M(\ce{O})=\qty{16}{\g\per\mol}$
          $M(\ce{S})=\qty{32}{\g\per\mol}$.
\end{enumerate}

